\documentclass[10pt]{article}
\usepackage[a4paper,margin=19mm]{geometry}
\usepackage[no-math]{fontspec}\setmainfont{Latin Modern Roman}
\usepackage{amsmath,amssymb,amsthm,mathrsfs,enumitem}
\usepackage[bookmarks=false]{hyperref}\hypersetup{hidelinks}
\newtheorem{thm}{Theorem}
\newtheorem{lem}[thm]{Lemma}
\newtheorem{prop}[thm]{Proposition}
\newtheorem{cor}[thm]{Corollary}
\theoremstyle{definition}\newtheorem{definition}[thm]{Definition}
\newtheorem{rmk}[thm]{Remark}
\numberwithin{equation}{section}
\newcommand{\R}{\mathbb{R}}\newcommand{\C}{\mathbb{C}}\newcommand{\N}{\mathbb{N}}
\newcommand{\h}{{\jepPubScript{H}}}\newcommand{\emb}{\hookrightarrow}\newcommand{\cemb}{\Subset}
\newcommand{\eps}{\varepsilon}\newcommand{\p}{\partial}\newcommand{\n}{\nabla}\newcommand{\la}{\langle}\newcommand{\ra}{\rangle}
\renewcommand{\qedsymbol}{$\square$}
\emergencystretch=3em
\makeatletter
\renewcommand\subsection{\@startsection{subsection}{2}{\z@}{-3.25ex\@plus -1ex\@minus -.2ex}{1.5ex\@plus .2ex}{\normalfont\normalsize\bfseries}}
\makeatother
\DeclareRobustCommand{\jepPubScript}[1]{\begingroup\def\jepArg{#1}\def\jepTestA{A}\def\jepTestE{E}\def\jepTestO{O}\def\jepTestY{Y}\def\jepTestH{H}\ifx\jepArg\jepTestA\mathscr{A}\else\ifx\jepArg\jepTestE\mathscr{E}\else\ifx\jepArg\jepTestO\mathscr{O}\else\ifx\jepArg\jepTestY\mathscr{Y}\else\ifx\jepArg\jepTestH\mathscr{H}\else\PackageError{published-font}{Unsupported literal}{Only A E O Y H are inventoried.}\fi\fi\fi\fi\fi\endgroup}
\newcommand{\jepPubBullet}{{\scriptscriptstyle\bullet}}
\begin{document}
\begin{center}{\Large Stable item 1910: checked published transcription}\end{center}
\paragraph{Source and attribution.} Lucas Ambrozio, Reto Buzano, Alessandro Carlotto and Ben Sharp, \emph{Bubbling analysis and geometric convergence results for free boundary minimal surfaces}, Journal de l'École polytechnique -- Mathématiques 6 (2019), 621--664, DOI 10.5802/jep.102. \href{https://jep.centre-mersenne.org/articles/10.5802/jep.102/}{Publisher article}; \href{https://creativecommons.org/licenses/by/4.0/}{CC BY 4.0}.
\paragraph{Selected independent problem.} Original Theorem 33, local extraction of a nontrivial bubble or half-bubble at an unstable boundary scale for $2\leq n\leq6$, with all original scale, connected-component, eigenvalue and stability hypotheses, both coordinate cases and the final nested-region assertion. Introductory definitions, complete Section 2, and Section 3 through the complete proof of Theorem 33 are retained. The complete original 45-page publisher PDF and bibliography are supplied. Surrounding global quantization and neck analysis are excluded from the selected claim.
\paragraph{Difficulty.} The boundary blow-up requires local exhaustion compactness, transfer of stability and negative eigenvalues between normal and Fermi coordinates, and exclusion of planar limits. After standard compactness and stable curvature estimates, the author still supplies the boundary extraction and nested unstable-region transfer argument. These source-local geometric and spectral limits exceed routine qualification-exam manipulations; the selection is H2.
\paragraph{Transcription notice.} Publisher-authority checked editable transcription with original numbering and formulas. Typography and page breaks are adapted using installed fonts. Supporting results are uncounted context. The nonexclusive arXiv offer is a private unexecuted aid and is excluded from public exports; it grants no native redistribution rights.
\section*{Original introductory definitions}
First of all, we recall some basic definitions and notation. Given a compact Riemannian manifold $(N^{n+1},g)$ of dimension $n+1\geqslant 3$, with non-empty boundary $\partial N$, let $M^n$ be a codimension one properly embedded submanifold: we say that $M^n$ is a free boundary minimal hypersurface in $(N^{n+1},g)$ if it is a critical point for the $n$-dimensional area functional under variations satisfying the sole constraint that $\partial M\subset\partial N$ or, equivalently, if $M$ has zero mean curvature and meets the ambient boundary orthogonally. In order to avoid ambiguities, let us remark that properness is always understood here in the strong sense that $M \cap \partial N=\partial M$ (which is consistent with \cite{FL14, LZ16B, ACS17, ACS18b} among others).
    We let $\mathfrak{M}$ denote the class of smooth, connected, and properly embedded free boundary minimal hypersurfaces and consider the subclasses given by
     	\[
     		\mathfrak{M}(\Lambda,I) := \{ M \in \mathfrak{M} \ : \ \h^n(M) \leqslant \Lambda,\ \operatorname{index}(M) \leqslant I\},
     			\]
     			so with the additional requirements that area and Morse index be bounded from above (cf. \cite{Sha15}); and for $p$ an integer greater or equal than one
     \[
     			\mathfrak{M}_p(\Lambda,\mu) := \{ M \in \mathfrak{M} \ : \ \h^n(M) \leqslant \Lambda,\ \lambda_p \geqslant -\mu\},
     			\]
     			so with area bounded from above and $p^{th}$ eigenvalue $\lambda_p$ of the Jacobi operator bounded from below (cf. \cite{ACS15}).	
     			
     			
     		
     		In \cite{ACS17}, Ambrozio-Carlotto-Sharp proved a compactness theorem for $\mathfrak{M}_p(\Lambda,\mu)$ when the ambient dimension is less than eight: given a sequence of free boundary minimal hypersurfaces $\left\{M_k\right\}$ in $\mathfrak{M}_p(\Lambda,\mu)$, there is some smooth limit to which the sequence sub-converges smoothly and graphically (with finite multiplicity $m$) away from a discrete set ${\jepPubScript{Y}}$ on the limit, where curvature concentration occurs. These results apply, as a special case, to all classes of the form $\mathfrak{M}(\Lambda,I)$ (since this is clearly the same as $\mathfrak{M}_{I+1}(\Lambda,0)$). 
     			A significant part of the present article is aimed at an accurate description of the local picture around those points of bad convergence, in analogy with the results obtained, for the closed case, by Buzano and Sharp in \cite{BS17}. We then specify these results to the three dimensional scenario, and derive various sorts of new geometric results based on these tools.	\\
     			
     			
     			
     		As a first step in this program, we quantify the lack of smooth compactness (via a blow-up argument) by a finite number of non-trivial, complete, properly embedded minimal hypersurfaces of finite total curvature $\Sigma^n\emb \R^{n+1}$ as in \cite{BS17}, except this time it is possible that only `half' of $\Sigma$ is captured at a boundary point of bad convergence. In order to formalize this picture, let us introduce some terminology:
     		we employ the word \textsl{bubble} to denote a complete, connected and properly embedded minimal hypersurface of finite total curvature in $\R^{n+1}$; we employ the word \textsl{half-bubble} to denote a complete, connected, properly embedded minimal hypersurface that is contained in a half-space and has (non-empty) free boundary with respect to the boundary of this half-space,  and has finite total curvature. In both cases, the total curvature is understood to be the integral of the $n$-th power of the length of the second fundamental form, the corresponding functional being denoted by ${\jepPubScript{A}}(\cdot)$. Furthermore, we shall say that a (half-)bubble is \textsl{non-trivial} if it is not flat, namely if it is not a (half-)hyperplane.
     		
     		
\setcounter{section}{1}\setcounter{thm}{10}
\section{Preliminary results on half-bubbles}\label{sec:prelim}
        
        
        	Let $\Sigma^n\subset \R^{n+1}$ be a half-bubble in the sense above. Let us denote by $\Pi_1$ the closed half-space bounded by $\Pi$ that contains the interior of $\Sigma$ and by $\Pi_2$ the other closed half-space bounded by $\Pi$. The terminology \textsl{half-bubble} can be justified as follows: if we reflect $\Sigma$ across $\Pi$ in $\R^{n+1}$ we get a minimal hypersurface without boundary, which a priori is only $C^1$, but a posteriori is $C^{1,\alpha}$ by means of a standard application of De Giorgi-Nash estimates, hence smooth by Schauder theory. Such a minimal hypersurface shall be denoted by $\check{\Sigma}$ and be referred to as the \textsl{double} of $\Sigma$.
        	
        	
        	In relation to the geometric applications we are about to present, we need to compare the Morse index of $\Sigma$ (as a free boundary minimal surface, thus with suitable boundary conditions) and the Morse index of $\check{\Sigma}$. In fact, our results will follow as a specification of a more general discussion.
        
     
        	
        	\subsection{Schr\"odinger-type operators on involutive manifolds}\label{subs:general}
        	
        	Let $(\check{\Sigma}^{n},g)$ be a complete Riemannian manifold, without boundary (in Section \ref{subs:morsehalf} we will then specify our discussion to the ambient manifold $\check{\Sigma}$ presented in the previous paragraph, namely a bubble with its induced metric). Suppose there exists a Riemannian involution $\tau:(\check{\Sigma},g) \to (\check{\Sigma},g)$, namely a smooth isometry satisfying the two conditions:
        	\[
        	\tau\circ\tau = \operatorname{id}, \ \ \tau\neq \operatorname{id}.
        	\] 
        	
        Given a linear functional space $X$ consisting of functions defined on $(\check{\Sigma},g)$ let us then introduce the subspaces of even and odd functions with respect to the action of $\tau$:
        	\[
        	X_{\jepPubScript{E}}=\left\{\varphi\in X: \varphi\circ\tau=\varphi \right\}, \ X_{\jepPubScript{O}}=\left\{\varphi\in X: \varphi\circ\tau=-\varphi \right\}.
        	\]	
        	In particular, notice that 
        	\[
        	C^{\infty}=C^{\infty}_{{\jepPubScript{E}}}\oplus C^{\infty}_{{\jepPubScript{O}}}.
        	\]
        	This direct sum is orthogonal in $L^2$ when restricted to the subspace of smooth, compactly supported functions. In fact, we have
        	\[
        	L^2=L^2_{{\jepPubScript{E}}}\oplus^{\perp} L^2_{{\jepPubScript{O}}}.
        	\]
        Given a function $V\in C^{\infty}_{{\jepPubScript{E}}}$ we wish to study a Schr\"odinger-type operator of the form
        	\[
        	T_V\varphi=\Delta_g\varphi+V\varphi, 
        	\]
        	together with the associated quadratic form
        	\[
        	Q_V(\varphi,\varphi)=-\int_{\check{\Sigma}}\varphi T_V \varphi\,d\h^n=\int_{\check{\Sigma}}(|\nabla \varphi|^2-V\varphi^2)\,d\h^n,
        	\]
        	which a priori is defined on the set of smooth, compactly supported functions.
        	
        	
        	\begin{definition}\label{def:index} In the setting above, we define
        		\[
        		\operatorname{Ind} (Q_V), \ \ \operatorname{Ind}_{{\jepPubScript{E}}} (Q_V), \ \ \operatorname{Ind}_{{\jepPubScript{O}}} (Q_V)
        		\]
        		as the largest dimension of a linear subspace of 
        		\[
        		C^{\infty}_c, \ \ (C^{\infty}_c)_{{\jepPubScript{E}}}, \ \ (C^{\infty}_c)_{{\jepPubScript{O}}}
        		\]
        		respectively, where the quadratic form $Q_V$ is negative definite.
        	\end{definition}	
        	
        	
        	Thereby, the following inequality is straightforward:
        	
        	\begin{lem}\label{lem:ineq}
        		In the setting above, 
        		\[
        		\operatorname{Ind} (Q_V) \geqslant \operatorname{Ind}_{{\jepPubScript{E}}} (Q_V)+\operatorname{Ind}_{{\jepPubScript{O}}} (Q_V)
        		\]
        		and
        		\[
        		\operatorname{Ind} (Q_V)=0 \iff \operatorname{Ind}_{{\jepPubScript{E}}} (Q_V)=0 \ \text{and} \ \operatorname{Ind}_{{\jepPubScript{O}}} (Q_V)=0.
        		\] 
        	\end{lem}
        	
        	\begin{proof} Since $V$ is even and $\tau$ is an isometry, it is immediate to check that the operator $T_V$ preserves the decomposition of $C^{\infty}$ into even and odd functions. Thus, if $\varphi_{{\jepPubScript{E}}}\in (C^{\infty}_c)_{{\jepPubScript{E}}}$ and $\varphi_{{\jepPubScript{O}}}\in (C^{\infty}_c)_{{\jepPubScript{O}}}$ then
        		\begin{equation}\label{eq:mixedprod}
        		Q_{V}(\varphi_{{\jepPubScript{E}}},\varphi_{{\jepPubScript{O}}})=0,
        		\end{equation}
        		and hence, by bilinearity, given any $\varphi\in C^{\infty}_c$ and writing $\varphi=\varphi_{{\jepPubScript{E}}}+\varphi_{{\jepPubScript{O}}}$ one has
        		\[
        		Q_{V}(\varphi,\varphi)=Q_{V}(\varphi_{{\jepPubScript{E}}},\varphi_{{\jepPubScript{E}}})+Q_{V}(\varphi_{{\jepPubScript{O}}},\varphi_{{\jepPubScript{O}}}),
        		\]
        		which easily implies the claims.
        	\end{proof}		
        	
        	
        	
        	In fact, we claim that the inequality above is actually an equality. 
        	With that goal in mind, we restrict to the case of \textsl{finite} Morse index and recall a simple but useful result:
        	
        	\begin{prop}\label{pro:l2spec}(cf. \cite[Prop.~2]{FC85})
        		In the setting above, the following two statements are equivalent:
        		\begin{enumerate}
        			\item{the index of the quadratic form $Q_V$ is finite;}
        			\item{there exists a finite dimensional subspace $W$ of $L^2$ having an orthonormal basis $\varphi^1,\ldots,\varphi^k$ consisting of eigenfunctions of $T_V$ with eigenvalues $\lambda_1,\ldots,\lambda_k$ respectively. Each $\lambda_i$ is negative and $Q_V(\varphi,\varphi)\geqslant 0$ whenever $\varphi\in C^{\infty}_c\cap W^{\perp}$. In this case $k$ equals the Morse index of $Q_V$.}	
        		\end{enumerate}		
        	\end{prop}
        	
        	
        	\begin{rmk}\label{rmk:Vbounded} If $V$ is assumed to be bounded, then the basis provided in Proposition \ref{pro:l2spec} actually consists of elements of finite Dirichlet energy, i.e. functions belonging to $W^{1,2}$, the closure of $C^{\infty}_c$ with respect to the Sobolev norm determined by
        		\begin{equation}\label{eq:SobNorm}
        		\|\varphi\|^2_{W^{1,2}}:= \|\varphi\|^2_{L^2}+\|\nabla \varphi\|^2_{L^2}.
        		\end{equation}	
        	\end{rmk}	
        	
        	
        	
        	Based on this fact and in view of the geometric applications we wish to present, we \emph{assume} from now onwards that the function	$V\in C^{\infty}_{{\jepPubScript{E}}}$ is uniformly bounded. 
        
        	
        	It is convenient to introduce the relevant notion of (point) spectrum in this setting.
        	
        	\begin{definition}\label{def:spec} In the setting above, we define
        		\[
        		\operatorname{spec} (Q_V):=\Big\{\text{critical values of the map} \ \varphi\in W^{1,2}\smallsetminus \left\{0\right\} \mapsto \frac{Q_V(\varphi,\varphi)}{\|\varphi\|^2_{L^2}}\Big\};
        		\]
        		\[
        		\operatorname{spec}_{{\jepPubScript{E}}} (Q_V):=\Big\{\text{critical values of the map} \ \varphi\in (W^{1,2})_{{\jepPubScript{E}}}\smallsetminus \left\{0\right\} \mapsto \frac{Q_V(\varphi,\varphi)}{\|\varphi\|^2_{L^2}}\Big\};
        		\]
        		\[
        		\operatorname{spec}_{{\jepPubScript{O}}} (Q_V):=\Big\{\text{critical values of the map} \ \varphi\in (W^{1,2})_{{\jepPubScript{O}}}\smallsetminus \left\{0\right\} \mapsto \frac{Q_V(\varphi,\varphi)}{\|\varphi\|^2_{L^2}}\Big\}.
        		\]
        	\end{definition}
        To avoid ambiguities: $\lambda\in \operatorname{spec}(Q_V)$ if and only if we can find an associated critical point $\varphi \in W^{1,2}\smallsetminus\left\{0\right\}$, which satisfies
        	\begin{equation}\label{eq:weak}
        	\int_{\check{\Sigma}}(\nabla \varphi \cdot \nabla\zeta -V\varphi\zeta)\,d\h^n = \lambda \int_{\check{\Sigma}}\varphi\zeta\,d\h^n, \ \ \ \forall \zeta\in C^{\infty}_c,
        	\end{equation}
        	whence it is standard to get that $\varphi$ is actually smooth and solves, in a classical sense, the eigenvalue equation
        	\[
        	\Delta_g \varphi+V\varphi+\lambda\varphi=0
        	\]
        	which also implies that $\Delta_g\varphi \in L^2$. 
        	
        	Similar arguments can be applied to the other two cases as well, with an important caveat. By definition, we have $\lambda\in \operatorname{spec}_{{\jepPubScript{E}}} (Q_V)$ (respectively $\lambda\in \operatorname{spec}_{{\jepPubScript{O}}} (Q_V)$) if equation \eqref{eq:weak} is satisfied for every \textsl{even} (respectively \textsl{odd}) test function $\zeta\in C^{\infty}_c$. However, we also have by symmetry arguments (cf. equation \eqref{eq:mixedprod}) that the same equation is satisfied for every \textsl{odd} (respectively \textsl{even}) test function, and hence for any $\zeta\in C^{\infty}_c$ (in either of the two cases). Therefore, we shall actually conclude
        	\[
        	\lambda\in \operatorname{spec}_{{\jepPubScript{E}}} (Q_V) \ \iff \ \exists \varphi\in  C^{\infty}_{{\jepPubScript{E}}}\smallsetminus \left\{0\right\} \ \text{such that} \	\Delta_g \varphi+V\varphi+\lambda\varphi=0,
        	\]
        	and
        	\[
        	\lambda\in \operatorname{spec}_{{\jepPubScript{O}}} (Q_V) \ \iff \ \exists \varphi\in  C^{\infty}_{{\jepPubScript{O}}}\smallsetminus \left\{0\right\} \ \text{such that} \	\Delta_g \varphi+V\varphi+\lambda\varphi=0.
        	\]
        	
        	
        	
        	
        	\begin{lem}\label{lem:specadd}
        		In the setting above,
        		\[
        		\operatorname{spec} (Q_V)=\operatorname{spec}_{{\jepPubScript{E}}} (Q_V) \cup 	\operatorname{spec}_{{\jepPubScript{O}}} (Q_V).
        		\]
        		Moreover, under the assumption that the operator $T_V$ has finite Morse index on $(\check{\Sigma},g)$, we have
        		\[
        		\operatorname{Ind} (Q_V) = \operatorname{Ind}_{{\jepPubScript{E}}} (Q_V)+\operatorname{Ind}_{{\jepPubScript{O}}} (Q_V).
        		\]
        	\end{lem}
        	
        		
        	
        	\begin{proof}
        		For the first claim, notice that the inclusion $\supseteq$ is obvious so let us discuss the other one. Let then $\lambda \in \operatorname{spec} (Q_V)$ and write $\varphi=\varphi_{{\jepPubScript{E}}}+\varphi_{{\jepPubScript{O}}}$ (the $L^2$ decomposition of an associated eigenfunction into even and odd parts), thus by linearity
        		\[
        		0=\Delta_g \varphi+V\varphi+\lambda\varphi= \underbrace{(\Delta_g \varphi_{{\jepPubScript{E}}}+V\varphi_{{\jepPubScript{E}}}+\lambda\varphi_{{\jepPubScript{E}}})}_{\in L^2_{{\jepPubScript{E}}} } + \underbrace{(\Delta_g \varphi_{{\jepPubScript{O}}}+V\varphi_{{\jepPubScript{O}}}+\lambda\varphi_{{\jepPubScript{O}}})}_{\in L^2_{{\jepPubScript{O}}}}
        		\]
        		hence each of the two summands must vanish and thus (since either $\varphi_{{\jepPubScript{E}}}\neq 0$ or $\varphi_{{\jepPubScript{O}}}\neq 0$) we get $\lambda\in \operatorname{spec}_{{\jepPubScript{E}}} (Q_V)$ or $\lambda\in \operatorname{spec}_{{\jepPubScript{O}}} (Q_V)$ as it was claimed.
        		
        		
        		
        		For the second part, recall that by Proposition \ref{pro:l2spec} there exist (under the assumption that $T_V$ has finite Morse index) finitely many (say $k$) eigenfunctions $\varphi^1,\ldots,\varphi^k \in L^2$ that correspond to the negative eigenvalues $\lambda_1,\ldots,\lambda_k$ (where it is understood that each eigenvalue can be repeated if it comes with multiplicity). 
        		Following the argument we just presented, replace each $\varphi^j$ by means of the couple $\varphi^j_{{\jepPubScript{E}}}, \varphi^j_{{\jepPubScript{O}}}$ and set
        		\[
        		V=\operatorname{span}_{\R}\left\{\varphi^1,\ldots,\varphi^k \right\}, \ \ \widetilde{V}=\operatorname{span}_{\R}\left\{\varphi^1_{{\jepPubScript{E}}}, \varphi^1_{{\jepPubScript{O}}},\ldots, \varphi^k_{{\jepPubScript{E}}}, \varphi^k_{{\jepPubScript{O}}} \right\}.
        		\] 
        		Now, since $\varphi^j=\varphi^j_{{\jepPubScript{E}}}+\varphi^j_{{\jepPubScript{O}}}$ for any $j\in\left\{1,\ldots, k\right\}$ we have
        		\[
        		\text{dim}_{\R}(\widetilde{V})\geqslant k
        		\]
        		and thus there are at least $k$ functions in the collection $\left\{\varphi^1_{{\jepPubScript{E}}}, \varphi^1_{{\jepPubScript{O}}},\ldots, \varphi^k_{{\jepPubScript{E}}}, \varphi^k_{{\jepPubScript{O}}} \right\}$ which are not zero, and eigenfunctions either for $\operatorname{spec}_{{\jepPubScript{E}}} (Q_V)$  or for $\operatorname{spec}_{{\jepPubScript{O}}} (Q_V)$, with negative eigenvalues. Thereby, the inequality one gets, combined together with Lemma \ref{lem:ineq}, enables us to complete the proof.
        	\end{proof}	
        	
        	
        	
        	\subsection{The Morse index of half-bubbles}\label{subs:morsehalf}
        	
        	The discussion presented in the previous section directly applies, as a special case, to the study of the Jacobi operator
        	that is associated to the symmetrized complete minimal hypersurface that is obtained by reflecting a half-bubble. 
        	
        	Given a half-bubble $\Sigma$ we define, in analogy with Definition \ref{def:index} above, its Morse index considering the Jacobi form
\[        	
Q_{|A|^2}(u,u):=\int_{\Sigma}(|\nabla u|^2-|A|^2u^2)\,d\h^n.
\] 
More precisely, we give the following:
        	
        	\begin{definition} In the setting above, we define
        		$\operatorname{index} (\Sigma)$ to be the largest dimension of a linear subspace of $C^{\infty}_c(\Sigma)$ where $Q_{|A|^2}$ is negative definite. 
        	\end{definition}	
        	\noindent Notice that here the support of $\varphi\in C^{\infty}_c(\Sigma)$ can intersect $\partial\Sigma$ so that we are imposing no condition along $\partial\Sigma$, and $\operatorname{index} (\Sigma)$ is the (standard) Morse index of $\Sigma$ as a free boundary minimal hypersurface.
        	
        	\begin{lem}\label{lem:II1}
        A two-dimensional half-bubble has finite Morse index and this equals $\operatorname{Ind}_{{\jepPubScript{E}}}(Q_{|\check{A}|^2})$ where $\check{\Sigma}$ is the double of $\Sigma$.
         The same conclusion holds true in all dimensions for any half-bubble $\Sigma$ having Euclidean volume growth, meaning that there exists a constant $C=C(\Sigma)$ such that $\h^n(\Sigma\cap B_R(0))\leqslant C R^n$ for all $R\geqslant 0$.
        	\end{lem}
        
           \begin{proof}
        	The argument for the first part ($n=2$) goes as follows: by definition $\Sigma$ has finite total curvature, so its double $\check{\Sigma}$ will also have finite total curvature, hence finite Morse index (again by \cite{FC85}); thus Lemma \ref{lem:ineq} implies that the corresponding \textsl{even} Morse index (that is to say: the Morse index on even test functions, as per Definition \ref{def:index}) is also finite. Now, going back to the definitions it is clear that by restriction $\operatorname{index} (\Sigma)\geqslant \operatorname{Ind}_{{\jepPubScript{E}}}(Q_{|\check{A}|^2})$.
        	
        	For the converse inequality, we argue by contradiction as follows. If it were $\operatorname{index} (\Sigma)> \operatorname{Ind}_{{\jepPubScript{E}}}(Q_{|\check{A}|^2})$, we could consider a basis $\left\{\varphi^1,\ldots, \varphi^I\right\}$ for a subspace of $C^{\infty}_c(\Sigma)$, having dimension equal to $I:=\operatorname{Ind}_{{\jepPubScript{E}}}(Q_{|\check{A}|^2})+1$, where $Q_{|A|^2}$ is negative definite. 
        			If we extend each function by even symmetry, we obtain a family of even functions $\left\{\check{\varphi}^1,\ldots \check{\varphi}^I\right\}$ where the Jacobi form of $\check{\Sigma}$ is negative definite; each function is smooth away from $\partial\Sigma$  and we can consider for each $i\in\left\{1,\ldots, I\right\}$ an approximating sequence of smooth functions $\left\{\check{\varphi}^i_k\right\}$ that also have compact support, and such that $\check{\varphi}^i_k\to\check{\varphi}^i$ uniformly as one lets $k\to\infty$.
        	 		 We claim that for any $k$ large enough the family $\left\{\check{\varphi}^1_k,\ldots \check{\varphi}^I_k\right\}$ is linearly independent, which would then violate the definition of $\operatorname{Ind}_{{\jepPubScript{E}}}(Q_{|\check{A}|^2})$ and conclude the proof.
        			 If the claim were false, we could write down (for every $k\in\mathbb{N}$) a linear equation of the form
        		 \[
        		 \sum_i a^i_k \check{\varphi}^i_k=0, \ \text{where} \ \sum_i (a^i_k)^2=1,
        		 \]
        		whence passing to the limit for $k\to\infty$, we end up finding a (non-trivial) linear relation involving 
        	$\left\{\varphi^1,\ldots, \varphi^I\right\}$, that is impossible since this family of functions was chosen to be a basis.
        	 In higher dimensions, one can follow the same argument modulo invoking the main theorem of J. Tysk in \cite{T89}.
        	\end{proof}
       
        
        	
        	We can instead consider the Morse index of $\Sigma$ with Dirichlet boundary conditions:
        	
        	\begin{definition} In the setting above, we define
        		$\operatorname{index}_{\jepPubBullet} (\Sigma)$ to be the largest dimension of a linear subspace of $C^{\infty}_c(\mathring{\Sigma})$ where $Q_{|A|^2}$ is negative definite.
        	\end{definition}
        	
        	Here $\mathring{\Sigma}:=\Sigma\smallsetminus\partial\Sigma$ is the interior of $\Sigma$, and so this is the standard notion of Morse index for minimal surfaces with respect to variations that fix the boundary. Following the same conceptual scheme as above, we have the following ancillary result:
        	
        	\begin{lem}\label{lem:II2}
        		A two-dimensional half-bubble has finite Morse index \emph{with Dirichlet boundary conditions} and this equals $\operatorname{Ind}_{{\jepPubScript{O}}}(Q_{|\check{A}|^2})$ where $\check{\Sigma}$ is the double of $\Sigma$.
        		The same conclusion holds true in all dimensions for any half-bubble $\Sigma$ having Euclidean volume growth, meaning that there exists a constant $C=C(\Sigma)$ such that $\h^n(\Sigma\cap B_R(0))\leqslant C R^n$ for all $R\geqslant 0$.
        	\end{lem}	
        	
        	\begin{proof}
        	Arguing as for the lemma above, let us focus on $n=2$. When $\Sigma$ is a two-dimensional half-bubble, the finiteness of $\operatorname{index}_{\jepPubBullet} (\Sigma)\leqslant \operatorname{index}(\Sigma)$ follows straight from Lemma \ref{lem:II1}. That being said, it is clear that we also have the inequality
        	$\operatorname{index}_{\jepPubBullet} (\Sigma)\leqslant \operatorname{Ind}_{{\jepPubScript{O}}}(Q_{|\check{A}|^2})$ by odd extension. The converse inequality
        	relies instead on the fact that odd functions must vanish along $\partial\Sigma=\operatorname{Fix}(\tau)$ so that one can consider a linearly independent set of $\operatorname{Ind}_{{\jepPubScript{O}}}(Q_{|\check{A}|^2})$ elements spanning a subspace on which $Q_{|\check{A}|^2}<0$, restrict each such function to $\Sigma$ and then truncate using a compactly supported cutoff function. We omit the standard details.
        		\end{proof}
        	
        Based on Lemma \ref{lem:II1} and \ref{lem:II2} one can rephrase the inequality $\operatorname{index}_{\jepPubBullet} (\Sigma)\leqslant \operatorname{index}(\Sigma)$ as
        	\begin{equation}\label{eq:EindexvsOindex}
        	\operatorname{Ind}_{{\jepPubScript{E}}}(Q_{|\check{A}|^2})\geqslant  \operatorname{Ind}_{{\jepPubScript{O}}}(Q_{|\check{A}|^2}),
        	\end{equation}
        	hence, using the second part of Lemma \ref{lem:specadd}, we get the estimate
        	\begin{equation}\label{eq:fundineq}
        	2 \operatorname{Ind}_{{\jepPubScript{O}}}(Q_{|\check{A}|^2})\leqslant \operatorname{Ind}(Q_{|\check{A}|^2})\leqslant 2 \operatorname{Ind}_{{\jepPubScript{E}}}(Q_{|\check{A}|^2}).
        	\end{equation}
        	From there, we can derive some interesting characterizations for the cases when the Morse index of $\check{\Sigma}$ (which is, in our notation, $	\operatorname{Ind}(Q_{|\check{A}|^2})$) equals 0, 1 or 2:
        	
        	\begin{cor}\label{cor:lowindex}
        		In the setting above, we have 
        		\begin{align*}
        		&\text{(1)}& 		\operatorname{Ind}(Q_{|\check{A}|^2})=0 \ \iff \operatorname{Ind}_{{\jepPubScript{E}}}(Q_{|\check{A}|^2})=0 \ \text{and} \ \operatorname{Ind}_{{\jepPubScript{O}}}(Q_{|\check{A}|^2})=0; \\
        		&\text{(2)}& 		\operatorname{Ind}(Q_{|\check{A}|^2})=1 \ \iff \operatorname{Ind}_{{\jepPubScript{E}}}(Q_{|\check{A}|^2})=1 \ \text{and} \ \operatorname{Ind}_{{\jepPubScript{O}}}(Q_{|\check{A}|^2})=0;\\
        		&\text{(3)}&		\operatorname{Ind}(Q_{|\check{A}|^2})=2 \ \iff 
        				\begin{cases}\operatorname{Ind}_{{\jepPubScript{E}}}(Q_{|\check{A}|^2})=2 \ \text{and} \ \operatorname{Ind}_{{\jepPubScript{O}}}(Q_{|\check{A}|^2})=0 \ \ \textbf{or} \\ 
        				 \operatorname{Ind}_{{\jepPubScript{E}}}(Q_{|\check{A}|^2})=1 \ \text{and} \ \operatorname{Ind}_{{\jepPubScript{O}}}(Q_{|\check{A}|^2})=1.
        				\end{cases}
        				\end{align*}
        	\end{cor}
        	
        	So far our discussion is applicable to hypersurfaces of dimension $n\geqslant 2$, and we shall now specify them to the special case $n=2$, where some interesting consequences can be drawn.
        	Indeed, we can turn classification results for complete minimal surfaces in $\R^{3}$ into classification results for free boundary minimal hypersurfaces contained in a half-space. The first one is a Bernstein-type theorem for half-bubbles:
        	
        	\begin{cor}\label{cor:stable}
        		A two-dimensional, stable half-bubble is isometric to a half-plane.	
        	\end{cor}	
        	
        	\begin{proof}
        		It suffices to combine (1) of Corollary \ref{cor:lowindex}, inequality \eqref{eq:EindexvsOindex} and the characterization of stable minimal surfaces in $\R^{3}$ (see \cite{dCP79,FCS80,Pog81}).
        	\end{proof}
        	
        	\begin{rmk}\label{rmk:higher-stable}
        	In fact, the stability estimates by Schoen-Simon allow to obtain a higher-dimensional counterpart of the previous result: \textsl{Let $\Sigma^n\subset\R^{n+1}$, $2\leqslant n\leqslant 6$ be a half-bubble. If $\Sigma$ is stable and has Euclidean volume growth, then $\Sigma$ is a hyperplane.}
        	\end{rmk}
        		
        	
        	
        	
        	Similarly, we get a simple result in the index one case:	
        	
        	\begin{cor}\label{cor:indexone}
        		A two-dimensional, index one half-bubble is isometric to a half-catenoid.	
        	\end{cor}	
        	
        	\begin{proof}
        		It suffices to combine inequality \eqref{eq:EindexvsOindex}, parts (2) and (3) of Corollary \ref{cor:lowindex}, the characterization of index one minimal surfaces in $\R^3$, see \cite{LR89}, and the non-existence result for index two complete minimal surfaces in $\R^3$, see \cite{CM14}.
        	\end{proof}
        	
        	In fact, inspired by recent work of Chodosh-Maximo \cite{CM14} we may wonder whether there exist two-dimensional half-bubbles whose Morse index equals two.
        	Obviously, a negative result would follow by proving that a half-bubble of Morse index 2 has vanishing Morse index with Dirichlet boundary conditions. Yet, this does not directly follow from the results above. 	
        	
        	\subsection{Geometry of half-bubbles}\label{subs:geomhalf}
        	
        	Let us discuss now some basic facts about the geometry and topology of half-bubbles. First of all, keeping in mind the quantization identity presented in Theorem 1, we wish to express the total curvature of a half-bubble in terms of its topology.
        	Clearly, in ambient dimension three ($n=2$) we can apply the Gauss equations to get
        	\begin{equation}\label{eq:half2nd}
        	2{\jepPubScript{A}}(\Sigma)={\jepPubScript{A}}(\check{\Sigma})=-2\int_{\check{\Sigma}} K_{\check{\Sigma}}\,d\h^2,
        	\end{equation}
        	while by Gauss-Bonnet (cf. Jorge-Meeks \cite{JM83}), denoting by $\gamma(\check{\Sigma})$ the genus of $\check{\Sigma}$ and by $b(\check{\Sigma})$ the number of its ends, we have
        	\begin{equation}\label{eq:GBcomplete}
        	\int_{\check{\Sigma}} K_{\check{\Sigma}}\,d\h^2=2\pi (\chi(\check{\Sigma})-b(\check{\Sigma}))=2\pi(2-2\gamma(\check{\Sigma})-2b(\check{\Sigma})) .
        	\end{equation}
        	It is then important to relate such data to the topological data of the half-bubble $\Sigma$. There are two easy examples of half-bubbles to be kept in mind throughout the following discussion, which we will prove in Lemma \ref{lem:euler} to be, roughly speaking, prototypical.
        	
        	\begin{definition}\label{def:horcurvercut} In the flat Euclidean space $\R^3$ we introduce the following terminology:
        		\begin{itemize} \item{the half-catenoid given by
        			\[
        			\left\{(x^1, x^2, x^3)\in \R^3 : \ (x^3)^2=\cosh ((x^1)^2+(x^2)^2), \ \ x^3\geqslant 0 \right\}
        			\]
        			will be called \textsl{horizontally cut half-catenoid};}	
       \item{the half-catenoid given by
        			\[
        		\left\{	(x^1, x^2, x^3)\in \R^3 : \ (x^3)^2=\cosh ((x^1)^2+(x^2)^2), \ \ x^1\geqslant 0 \right\}	
        			\]
        			will be called \textsl{vertically cut half-catenoid}.}
        			\end{itemize}
        	\end{definition}
        	
        	 Before proceeding further, it is helpful to recall some information about the asymptotic structure of complete minimal surfaces in $\R^3$.
        	
        	\begin{rmk}\label{rem:reginf}
        		Let $S$ be a complete, embedded minimal surface of $\R^3$. Then:
        		\begin{itemize}
        			\item{$S$ has finite total curvature if and only if it has finite Morse index (cf. \cite{FC85});}
        			\item{if $S$ satisfies either of these two equivalent assumptions then it is regular at infinity (cf. \cite{Sc83}), namely it can be decomposed, outside any sufficiently large compact set, into a finite number of connected components (named \textsl{ends}) and each such end can be described as a graph (over a plane in $\R^3$) of a defining function having an expansion of the form
        				\[
        				u(x')=a\log|x'|+b+\frac{c_1 x^1}{|x|^2}+\frac{c_2 x^2}{|x|^2}+O(|x'|^{-2}), \ |x'|\longrightarrow\infty
        				\]
        				in suitable coordinates $x=(x^1,x^2,x^3)$, where $x'=(x^1,x^2)$.
        			}
        		\end{itemize}
        		Notice that if $\Sigma$ is a half-bubble then both facts apply to its double $\check{\Sigma}$ and hence, a posteriori, one can talk about the ends of a half-bubble and this notion is well-defined. Furthermore, each end of a half-bubble will have a precise asymptotic description of the type above.
        	\end{rmk}
        	
        		\begin{rmk}
        			We can somewhat strengthen Lemma \ref{lem:II1} and extend to half-bubbles the aforementioned equivalence result by Fischer-Colbrie:
        			\end{rmk}
        		\begin{prop}\label{pro:equiv}
        			A two-dimensional half-bubble has finite index if and only if it has finite total curvature.	The same conclusion holds true in all dimensions provided Euclidean volume growth is assumed.
        		\end{prop}
        		One of the two implications has been proved above in Lemma \ref{lem:II1}, while the other follows by combining Lemma \ref{lem:specadd} and the inequalities in \eqref{eq:fundineq}.
        	
        	\
        	
        	Let us now proceed in our discussion, relating all half-bubbles to one of the two models given in Definition \ref{def:horcurvercut}.
        	
        	\begin{lem}\label{lem:euler}
        		Let $\Sigma$ be a half-bubble and let $\check{\Sigma}$ denote its double in $\R^3$. Then one of the following two alternative cases occurs:
        		\begin{enumerate}
        			\item{the number of ends of $\check{\Sigma}$ equals double the number of ends of $\Sigma$, and their Euler characteristics are related by the equation
        				\[
        				\chi(\check{\Sigma})=2\chi(\Sigma);
        				\]	
        			}
        			\item{the number of ends of $\check{\Sigma}$ equals the number of ends of $\Sigma$,  and their Euler characteristics are related by the equation
        				\[
        				\chi(\check{\Sigma})=2\chi(\Sigma)-b(\Sigma).
        				\]
        			}	
        			\end{enumerate}
        			In either case, one has that $\chi(\check{\Sigma})-b(\check{\Sigma})=2(\chi(\Sigma)-b(\Sigma))$.
        	\end{lem}	
        	
        	Clearly, both cases do occur: (1) is exemplified by the horizontally cut half-catenoid, while (2) is exemplified by the vertically cut half-catenoid. 
        	
        	\begin{proof}
        		Based on Remark \ref{rem:reginf} applied to the symmetrized bubble $\check{\Sigma}$ we have that either each of its ends is contained in one of the half-spaces $\Pi_1$ and $\Pi_2$ (provided we remove from $\check{\Sigma}$ a sufficiently large ball centered at the origin) or instead all of them intersect $\Pi$. By symmetry, it is clear that in the first case the number of ends of $\check{\Sigma}$ equals double the number of ends of $\Sigma$ and we can justify the equation for the Euler characteristic using the Gauss-Bonnet theorem: since $\Sigma$ is assumed to have finite total curvature we can write for $\Sigma_r:=\Sigma \cap \left\{|x|\leqslant r\right\}$
        		\[
        		\int_{\Sigma}K_{\Sigma}\,d\h^2=\lim_{r\to\infty}\int_{\Sigma_r} K_{\Sigma}\,d\h^2
        		\]
        		and we can compute the limit on the right-hand side along a sequence of generic radii $\left\{r_i\right\}$, so that the intersection $\Sigma\cap \left\{|x|=r_i\right\}$ is transverse and thus consisting of finitely many, smooth closed curves, hence
        		\[
        		\int_{\Sigma_{r_i}} K_{\Sigma}\,d\h^2=2\pi \chi(\Sigma)-2\pi b(\Sigma)+o(1)
        		\]
        		(where we have used the fact that the geodesic curvature of $\partial\Sigma\subset\Sigma$ is zero at all points).
        		Thus the conclusion comes straight from combining this equation with \eqref{eq:half2nd} and \eqref{eq:GBcomplete}. This is the scenario corresponding to case (1).
        		
        		If instead all ends of $\check{\Sigma}$ intersect the plane $\Pi$ then we are in case (2) and the claim concerning $\chi(\check{\Sigma})$ comes, once again, by an elementary computation: this time we have, along a sequence of generic radii 
        		\[
        		\int_{\Sigma_{r_i}} K_{\Sigma}\,d\h^2=2\pi \chi(\Sigma)-\pi b(\Sigma)-\pi b(\Sigma)+o(1)
        		\]
        		by keeping the contributions of both the geodesic curvature and of the exterior angles (at the non-smooth points of $\Sigma_r$) into account. In that respect, notice that  there are exactly $2b(\Sigma)$ exterior angles, each of them being $\pi/2+o(1)$ as $i\to\infty$. Thereby, the conclusion follows.
        	\end{proof}	
        	
        	
        	
        		
        		
        		
        		
        	
        
        	
        	
        	We can extend to half-bubbles some further classification results that are well-known for complete minimal surfaces in $\R^3$. This discussion parallels the one presented in the previous subsection, which was based on index-theoretic criteria instead.
        	
        	
        	\begin{cor}\label{cor:oneend}
        		A two-dimensional half-bubble with one end is isometric to either a half-plane or to a horizontally cut half-catenoid.
        	\end{cor}	
        	
        	\begin{proof}
        		Denoted by $\Sigma$ the half-bubble in question and $\check{\Sigma}$ its double, Lemma \ref{lem:euler} implies that \textsl{either} $\Sigma$ and $\check{\Sigma}$ have the same number of ends (one) \textsl{or} $\check{\Sigma}$ will have two ends. By classical results of Schoen (see \cite{Sc83}) the first case happens if and only if $\check{\Sigma}$ is a plane, and the second if $\check{\Sigma}$ is a catenoid. The conclusion in the former case is straightforward, so let us discuss the latter. The only way $\Sigma$ has exactly one end (which we are assuming) is that it is cut with a plane that is parallel to its two ends. But then the only way such a cut gives rise to a free boundary minimal surface is when the plane in question passes through the center of the catenoid, which proves the claim.
        	\end{proof}	
        	
        	\begin{cor}\label{cor:genuszero}
        		A two-dimensional half-bubble whose Euler characteristic equals one is isometric to either a half-plane or to a vertically cut half-catenoid.
        	\end{cor}	
        	
        	\begin{proof}
        		Like we did above, we need to consider two cases and apply the equations provided by Lemma \ref{lem:euler} for the Euler characteristics. If we are in case (1), the symmetrized bubble $\check{\Sigma}$ satisfies $\chi(\check{\Sigma})=2$ and this is not possible because the standard Jorge-Meeks formula for the Euler characteristic (cf. the second equality in \eqref{eq:GBcomplete}) forces $\chi(\check{\Sigma})\leqslant 1$. 
        		Instead, in case (2) we would get that the bubble $\check{\Sigma}$ has genus zero, which allows to invoke \cite{LR91}: $\check{\Sigma}$ is either a plane or a catenoid. At this stage, the conclusion comes at once by considering all possible planes of symmetry of a catenoid.
        	\end{proof}	
        	
        	
        	 	
        	 	
\section{Bubbling analysis}\label{sec:bubbling}
        	 	
        	 	
        	 	In this section, we prove Theorem 5 and in Section 4 we complete the proof of Theorem 1. Before we begin either of these proofs notice that given $y\in {\jepPubScript{Y}}\smallsetminus \partial N$ the bubbling and neck analysis in \cite{BS17} goes through precisely as before since it is of a local nature. In particular the conclusions of Theorems 1 and 5 hold over smooth domains $U\cemb N\smallsetminus\partial N$ such that ${\jepPubScript{Y}}\cap \partial U =\varnothing$ and $y\in {\jepPubScript{Y}}\smallsetminus \partial N$ respectively, which implies that the proofs of both of these theorems is complete in the case that ${\jepPubScript{Y}}\cap \partial N$ is empty. We need only focus on the setting $y\in {\jepPubScript{Y}}\cap \partial N$ from now on. 
        	 	
        	 	Our main result here is a blow-up theorem (Theorem \ref{thm:blowup} below) which will detect a non-trivial bubble or half-bubble in all regions of coalescing index at $\partial N$. 
        	 	This result is based on a localized version of the compactness result for free boundary minimal hypersurfaces from \cite{ACS17}. Once we have proved Theorem \ref{thm:blowup}, we can then follow a scheme related to the bubbling analysis for closed minimal hypersurfaces developed in \cite[\S3]{BS17} to construct various point-scale sequences which will detect \emph{all} the non-trivial bubbles and half-bubbles that develop at points of curvature concentration on $\partial N$, yielding Theorem 5. 
        	 	
        	 	In order to obtain the quantization result in Theorem 1 from our proof of Theorem 5, it remains to show that no curvature is lost in \emph{annular neck regions} between the bubble scales. This last step will be carried out in Section 4.\\
        	 	
        	 	
        	 	
        	 	In order to state our results in a unified manner, we will use the following notation and conventions throughout. Consistently with Section \ref{sec:prelim} denote by $\Pi_1(a):=\{x^1\geqslant a\}$ a closed half-space of $\R^{n+1}$ \emph{containing the origin} and by $\Pi(a)$ its boundary (though we drop the dependence on $a$ when it is irrelevant).  We will consider %
        	 	a relatively open sub-domain $C^{\R^{n+1}}_\varrho(a)\subset \Pi_1$,  %
        	 	defined by $C^{\R^{n+1}}_\varrho(a)=[a,a+\varrho)\times B^{\R^n}_\varrho (0)$. We will allow $\varrho=\infty$ in the sequel at which point the domain is simply $\Pi_1(a)$.  %
        	 	Now equip $C^{\R^{n+1}}_\varrho$ with a smooth metric $g$ such that it is Euclidean at $(a,0,\dots,0)$ and for all $z\in C^{\R^{n+1}}_\varrho\cap \Pi$, $\partial_{x^1}(z)$ is the inward unit normal to $\Pi$ with respect to $g$. We say that $(C^{\R^{n+1}}_\varrho,g)$ is an \emph{adapted} domain and in the sequel we will allow the parameters $\varrho$ and $g$ to vary under these constraints, and $a$ will largely remain fixed.   %
        	 	
        	 	
        	 	Such domains $(C^{\R^{n+1}}_\varrho,g)$ can be used to analyze local properties of free boundary minimal hypersurfaces in arbitrary $N$ without loss of generality via the use of
        	 	a \emph{Fermi-coordinate} neighborhood at the boundary of $N$: 
        	 	Let $-\nu$ be the inward pointing normal to $\partial N$ and $\gamma_1(N)>0$ be so that 
        	 	\begin{eqnarray*}
        	 		& F:[0,\gamma)\times \partial N \longrightarrow N&\\
        	 		&(t,z)\longmapsto \operatorname{Exp}^N_z (-t\nu)&
        	 	\end{eqnarray*} 
        	 	is injective for $t\in [0,\gamma_1)$. 
        	 	
        	 	Consider a normal coordinate neighborhood of $q$ in $\partial N$, $\{x_i\}_{i=2}^{n+1}$. For ease of notation, we will now choose $\gamma_0(N) >0$ to be smaller than both $\gamma_1$ defined above, and the injectivity radius of $\partial N$ so that, first of all, these boundary normal coordinates are defined on $B^{\partial N}_\gamma (q)$ for any $q\in \partial N$ and $\gamma \leqslant \gamma_0$. Second of all, for all $q\in \partial N$ and $\gamma\leqslant \gamma_0$, we may now extend these to a Fermi-coordinate neighborhood $C_{\gamma}(q) \cong [a,a+\gamma)\times B^{\R^n}_\gamma (0)$ via the exponential map giving $x^{1} = a+ \operatorname{Exp}^N_{(0,x^2,\dots,x^{n+1})} (-t\nu)$ for $0\leqslant t<\gamma$. 
        	 	
        	 	Notice that the metric $g$ in these coordinates coincides with the Euclidean metric at $(a,0,\dots,0)$. Moreover $\partial_{x^1}(z) = -\nu(z)$, in particular $g_{1j}(z)=0$, for all $z\in \partial N = \{x^1=a\}$ and $j\geqslant 2$. In particular $C_\gamma(q)\cong C^{\R^{n+1}}_\gamma(a)$ is an adapted domain.  \\
        	 	
        	 	For $V\subset (C^{\R^{n+1}}_\varrho,g)$ consider the set $\mathfrak{M}^{V}$ (which shall depend on both $C^{\R^{n+1}}_\varrho$ and the background metric $g$) of smooth, connected and properly embedded minimal hypersurfaces $P\subset V$, furthermore requiring that if $P\cap \Pi \neq \varnothing$ then $P$ is free boundary with respect to $\Pi$. At the level of regularity, we always tacitly assume $V$ to be the intersection of a \emph{smooth} domain in $\R^{n+1}$ with $\Pi_1$. Any variational properties of $P$ are computed with respect to compactly supported variations in $V$ -- i.e. free boundary variations on $V\cap \Pi$ as well as Dirichlet conditions on $\partial V\smallsetminus \Pi$. Following our introduction, we define
        	 	\begin{equation*}
        	 	\mathfrak{M}^{V}_p(\Lambda,\mu) = \big\{ P \in \mathfrak{M}^{V} :  \forall x\in C^{\R^{n+1}}_\varrho, R>0,  \ \h^n(P\cap B^{\R^{n+1}}_R(x))\leqslant \Lambda R^n, \ \text{ and }\lambda_p(P) \geqslant -\mu \big\}.
        	 	\end{equation*}
        	 	
        	 	
        	 	
        	 
        	 	We recall from \cite[Th.~29]{ACS17} that if $2\leqslant n\leqslant 6$ and $\{M_k\}\subset \mathfrak{M}_p(\Lambda, \mu) \subset \mathfrak{M}$ then there exists a smooth, connected, compact embedded minimal hypersurface $M\subset N$ meeting $\partial N$ orthogonally along $\partial M$, $m\in \mathbb{N}$ and a finite set ${\jepPubScript{Y}}\subset M$ with cardinality $|{\jepPubScript{Y}}|\leqslant p-1$  such that, up to subsequence, $M_k \to M$ locally smoothly and graphically on $M\smallsetminus {\jepPubScript{Y}}$ with multiplicity $m$. To avoid ambiguities, let us remark that from now on we will always assume that  ${\jepPubScript{Y}}$ is the \emph{minimal} such set, so that in particular the convergence is never smooth about any $y\in{\jepPubScript{Y}}$.
        	 	
        	 	We will require the following local version of that compactness theorem.
        	 	
        	 	\begin{lem}\label{lemma:local}
        	 		Let $2\leqslant n \leqslant 6$ and $\{(C_{k}, g_k)\}$ be a sequence of adapted domains where we set $C_k=C^{\R^{n+1}}_{\varrho_k}(a)$ with $\{\varrho_k\}$ monotonically increasing, and choose $\eps >0$ small enough that $B^{\R^{n+1}}_\eps(0)\cemb C_1$. We will assume that $g_k\to g$ smoothly for some limit metric $g$ on any relatively open subset $V\cemb C_k$. Suppose we have a sequence $P_k\in \mathfrak{M}^{C_k}_p(\Lambda,\mu)$ (for some fixed constants $\Lambda, \mu\in \R_{>0}$ and a positive integer $p$ independent of $k$) such that $P_k\cap B^{\R^{n+1}}_{\eps}(0)\neq \varnothing$.
        	 		Then, up to subsequence: 
        	 		
        	 		
        	 		\begin{enumerate}
        	 			\item For any relatively open $V$ such that $B_\eps^{\R^{n+1}}(0)\subset V\cemb C_k$ for some $k$, there exists a smooth, connected, and embedded minimal hypersurface $P\subset V$ where $P_k \to P$ locally smoothly and graphically, with multiplicity $m\in \N$, for all $x\in P\smallsetminus {\jepPubScript{Y}}$ where ${\jepPubScript{Y}}\subset P$ is a finite set with cardinality $|{\jepPubScript{Y}}|\leqslant p-1$. 
        	 			
        	 			 If $\partial P\neq\varnothing$, then $P$ meets $\Pi$ orthogonally along $\partial P$ and if $P \in \mathfrak{M}^V$, then $P \in \mathfrak{M}^V_p(\Lambda,\mu)$.
        	 			 
        	 			  Finally ${\jepPubScript{Y}}\neq \varnothing$ if and only if there exists $\{y_k\}\subset V$ with $y_k\to y\in P$, and $r_k\to 0$ so that $\operatorname{index}(P_k\cap B^{\R^{n+1}}_{r_k}(y_k))\geqslant 1$. In either case we say that $y\in {\jepPubScript{Y}}$. 
        	 			
        	 			
        	 			
        	 			\item Assuming that $\varrho_k \to \infty$, then there exists a limit $M$ which is a smooth, embedded  minimal hypersurface in $(\Pi_1,g)$.   
        	 		If we additionally assume $\liminf_{k\to\infty}\lambda_p(P_k)\geqslant 0$ then  $P \in \mathfrak{M}^{\Pi_1}$ implies $P \in \mathfrak{M}^{\Pi_1}(\Lambda,p-1)$.
        	 			
        	 			Further assuming $g=g_0$ is the Euclidean metric then $P$ is either a bubble or a half-bubble; if $P$ is not properly embedded, then $P\equiv \Pi$. 
        	 			
        	 			Finally if ${\jepPubScript{Y}}\neq \varnothing$ then $P$ must be a plane in $\Pi_1$ or a half-plane orthogonal to $\Pi$. 
        	 			
        	 		\end{enumerate}
        	 		
        	 	\end{lem}
        	 	
        	 	\begin{proof}
        	 		Part (1):        		 All the steps in the proof of \cite[Th.~29]{ACS17} are local, hence can be adopted almost verbatim with only some cosmetic changes to conclude all but the final statement concerning the equivalent characterization of the condition ${\jepPubScript{Y}}\neq\varnothing$. To prove that assertion, let us first observe that since $P$ is smooth and $V$ is compact we can pick $r>0$ so that $P\cap B^{\R^{n+1}}_r(z)\cap V$ is \emph{strictly} stable for all $z\in P$. Thus, if ${\jepPubScript{Y}}=\varnothing$ then the convergence would be smooth everywhere, hence in particular for $k$ large enough each hypersurface $P_k$ would be stable in all balls of radius $r/2$ centered at points in $V$ (as defined above) by smooth convergence. Hence no sequence as in the statement can actually exist. Instead, for the converse implication one just needs to invoke the interior and free boundary curvature estimates of Schoen-Simon \cite{SS81} (cf. \cite[Th.~19]{ACS17}).  
        	 		
        	 		\
        	 		
        	 		Part (2):
        	 		We can apply Part (1) to some compact exhaustion $\{V_{\ell}\}$ of $\Pi_1$ to conclude the first statement via a diagonal sequence argument.%
        	 		
        	 		If $P$ is proper, i.e. $P \in \mathfrak{M}^{\Pi_1}$, then by Part (1) $P\in \mathfrak{M}^V_p(\Lambda, \mu)$ for \emph{all} $V \subset \Pi_1$. We now check that in this case in fact $\operatorname{index}(P)\leqslant p-1$. Assuming the contrary guarantees the existence of a set $V$ with $\lambda_p(P\cap V)=\alpha<0$ and therefore, following the argument in \cite[p.~2599--2600]{ACS15}, we find that eventually $\lambda_p(P_k\cap V)\leqslant \alpha/2$, contradicting the assumption that $\liminf_{k\to\infty}\lambda_p(P_k)\geqslant 0$.  
        	 		
        	 		If $g=g_0$ is the Euclidean metric, then we distinguish two cases. If $P$ is proper, then exploiting the assumption of Euclidean volume growth, we conclude by Proposition \ref{pro:equiv} that it has finite total curvature. If on the other hand $P$ is not proper, i.e. an interior point of $P$ lies in $\Pi$, then by a maximum principle argument $P=\Pi$, and therefore $P$ trivially has finite total curvature. 
        	 		
        	 		Finally, if ${\jepPubScript{Y}}\neq \varnothing$ then we may as well assume that $P$ is properly embedded (if not then $P=\Pi$ and we are done). Since $P$ is properly embedded and $\Pi_1$ is simply connected $P$ is two-sided. When we couple this with the lack of smooth convergence, we conclude that the convergence must be of multiplicity $m\geqslant 2$ (via Allard's interior and boundary regularity cf. \cite[Ths.~5, 17]{ACS17}) which allows for the construction of a positive Jacobi field over compact subsets $W\cemb \Pi_1$ (not necessarily vanishing at the boundary of $W$). In turn, this implies that $\lambda_1(P\cap W)\geqslant 0$ and the limit is stable over all such $W$ (here we use the standard argument to establish the positiveness of a first eigenfunction, noting that both an interior and boundary Hopf maximum principle will be required); hence $P$ is stable in $\Pi_1$. The conclusion follows by the usual classification of stable hypersurfaces with Euclidean volume growth when $P$ has no boundary \cite{SS81}, and Remark \ref{rmk:higher-stable} when $P$ has a free boundary on $\Pi$. 
        	 		\end{proof}
        	 	
        	 	
        	 	We can now establish the first main result of this section as a consequence of the above statement as well as the localized compactness theory developed in \cite[\S2]{BS17} for the case of closed hypersurfaces. 
        	 	
        	 	 We have already seen in Lemma \ref{lemma:local} that bad points of convergence $y\in {\jepPubScript{Y}}$ correspond to coalescing regions of index, so these are the points that we analyze in detail here. We will denote by $B_r(p)$ a normal coordinate neighborhood of $p\in N$ and by $C_r(q)$ a Fermi-coordinate neighborhood of $q\in \partial N$. 
        	 
        	 	
        	 	\begin{thm}\label{thm:blowup}
        	 		Let $\{M_k\}\subset \mathfrak{M}_p(\Lambda, \mu) \subset \mathfrak{M}$ so that (up to subsequence) $M_k\to M$ for $M$ as in \cite[Th.~29]{ACS17} and choose $\delta>0$ so that 
        	 		\begin{equation*}
        	 		2\delta<\min\Big\{\inf_{ y_i\neq y_j\in {\jepPubScript{Y}}} \operatorname{dist}_g(y_i,y_j),\,\,\, \operatorname{inj}_N,\,\,\, \gamma_0\Big\}.
        	 		\end{equation*} Assume the existence of sequences $\{\varrho_k\},\{r_k\}\subset \R_{>0}$ satisfying $r_k\to 0$, $\varrho_k\leqslant \delta$, $\varrho_k/r_k\to \infty$ and $\{p_k\in M_k\}$ with $p_k\to y\in \partial N$. We assume furthermore that $P_k$ is some connected component of $M_k\cap B_{\varrho_k}(p_k)$ and it satisfies
        	 		\begin{itemize}
        	 			\item $P_k\cap B_{r_k + (r_k)^2}(p_k)$ is unstable for all $k$,
        	 			\item $P_k\cap B_{r_k/2}(z)$ is stable for all $z\in P_k\cap B_{\varrho_k}(p_k)$. 
        	 		\end{itemize}
        	 		
        	 		After possibly passing to a further subsequence, we are in one of the following cases: \\
        	 		
        	 		\framebox[1.8cm]{Case I:} $\operatorname{dist}_g(p_k,\partial N)/r_k\to\infty$. Choose normal coordinates for $N$ centered at $p_k$ on the geodesic balls $B_{\varrho_k}(p_k)$ and consider the rescaled hypersurfaces $\widetilde{P}_k \subset B^{\R^{n+1}}_{\varrho_k/r_k}(0)$ defined by
        	 		\begin{equation*}
        	 		\widetilde{P}_k  : = \frac{1}{r_k} \big(P_k\cap B_{\varrho_k}(p_k)- p_k\big).
        	 		\end{equation*}		 Then $\widetilde{P}_k$ must converge smoothly on any compact set with multiplicity one to a non-trivial bubble $\Sigma\subset \R^{n+1}$ with $\Sigma\cap B^{\R^{n+1}}_{3/2}(0)$ unstable. \\
        	 		
        	 		
        	 		
        	 		\framebox[1.8cm]{Case II:} $\operatorname{dist}_g(p_k,\partial N)/r_k\to c\geqslant 0$. Let $q_k\in\partial N$ be so that $\operatorname{dist}_g(q_k,p_k)=\operatorname{dist}_g(p_k,\partial N)$ and $s_k$ a point on the geodesic connecting $q_k$ to $p_k$ which is a distance $cr_k$ from $q_k$. Choose a Fermi-coordinate neighborhood (so that $a=-c$) $C_{\varrho_k}(q_k)$ and consider the rescaled hypersurfaces $\widetilde{P}_k \subset C^{\R^{n+1}}_{\varrho_k/r_k}(-c)$ defined by
        	 		\begin{equation*}
        	 		\widetilde{P}_k  : = \frac{1}{r_k} \big(P_k\cap C_{\varrho_k}(q_k)- s_k\big).
        	 		\end{equation*}   In this case $\widetilde{P}_k$ must converge smoothly on any compact set with multiplicity one to \emph{either a full or half-bubble} $\Sigma\subset \Pi_1(-c)$, non-trivial in both cases, with $\Sigma\cap B^{\R^{n+1}}_{3/2}(0)$ unstable. %
        	 		In the case that $n=2$, $\Sigma$ must be a non-trivial half-bubble. \\ 	
        	 		Moreover, in each of the above cases, we must have $\lambda_1(P_k \cap (B_{2r_k}(p_k)))\to -\infty$. \\
        	 		
        	 		
        	 		Finally, suppose that there is another point-scale sequence $(\widehat{p}_k, \widehat{r}_k)$ with $\widehat{r}_k \geqslant r_k$ for all $k$ and such that
        	 		\begin{itemize}
        	 			\item $P_k\cap (B_{\widehat{r}_k+ (\widehat{r}_k)^2}(\widehat{p}_k)\smallsetminus B_{2r_k}(p_k))$ is unstable for all $k$,
        	 			\item there exists some $C<\infty$ with $B_{\widehat{r}_k}(\widehat{p}_k)\subset B_{Cr_k}(p_k)$ for all $k$,
        	 		\end{itemize}
        	 		then we also have $\lambda_1(P_k\cap (B_{2\widehat{r}_k}(\widehat{p}_k)\smallsetminus B_{2r_k}(p_k))) \to -\infty$.
        	 	\end{thm}
        	 	
        	 	\begin{rmk}\label{rmk:coord}
        	 		Suppose that we are in the setting of Case II, and consider $B_{\varrho_k}(p_k)\cap C_{\varrho_k}(q_k)$. Letting $\Phi_k : B_{\varrho_k}(p_k) \to B^{\R^{n+1}}_{\varrho_k/r_k}(0)$ and $\Psi_k : C_{\varrho_k}(q_k)\to C^{\R^{n+1}}_{\varrho_k/r_k}(-c)$ denote the blown-up normal and Fermi-coordinates respectively, one can easily check that 
        	 		$$\Phi_k\circ \Psi_k^{-1}:  \Psi_k(B_{\varrho_k}(p_k)\cap C_{\varrho_k}(q_k)) \longrightarrow \Phi_k(B_{\varrho_k}(p_k)\cap C_{\varrho_k}(q_k))$$
        	 		converges locally smoothly on $\Pi_1$ to the identity as $k\to \infty$. Therefore the resulting blow-up and conclusions in Case II can actually be taken exactly as in Case I with respect to normal coordinates centered at $p_k$, or indeed as stated in the Theorem. 
        	 	\end{rmk}
        	 	
        	 	\begin{rmk}\label{rmk:equivballs}
        	 			Referring to the statement above, we further note that in these coordinates geodesic balls in $N$, $B_r(p)$, are directly comparable to Euclidean balls: for any $T>0$ there exists a sequence $\left\{\beta_k(T)\right\}$ with $\beta_k(T) \searrow 1$ so that for all $p\in B_{Tr_k}(s_k)$, if we denote $\widetilde{p}\in C^{\R^{n+1}}_{T}(-c)$ the point corresponding to $p$ in these rescaled coordinates then, by abusing notation,  
        	 			$$B^{\R^{n+1}}_{r\beta_k^{-1}}(\widetilde{p}) \subset B_r(p)\subset B^{\R^{n+1}}_{r \beta_k}(\widetilde{p}).$$
        	 		 This equivalence will be exploited along the course of the following proof when transferring certain variational properties (typically: eigenvalue bounds) back and forth between geodesic balls and coordinate balls. As it is well-known, the same equivalence holds true in geodesic normal coordinates as well.
        	 	\end{rmk}	
        	 	
        
        	 	
        	 	
        	 	
        	 	
        	 	
        	 	
        	 	
        	 	
        	 	\begin{proof}
        	 	
        	 		Case I: We can follow the arguments exactly as in \cite[Cor.~2.6]{BS17} to conclude all the statements of the theorem. \\
        	 		
        	 		
        	 		Case II: 
        	 		In these Fermi-coordinates, we are now in a position to apply Lemma \ref{lemma:local} and it remains to check that the limit $P$ is in fact a non-trivial (namely: it is not flat) and that the final statement of the theorem, concerning the point-scale sequence $(\widehat{p}_k,\widehat{r}_k)$, also holds. 
        	 		
        	 		
        	 		Concerning the first assertion, observe that by Lemma \ref{lemma:local} it suffices to check that $P$ is not stable. 
        	 		We have ${\jepPubScript{Y}}=\varnothing$ since $\widetilde{P}_k$ is stable on all balls of radius $1/4$ (by virtue of our assumption), hence $\widetilde{P}_k$ converges smoothly and graphically on every compact subset of $\Pi_1$ to a connected minimal hypersurface $P$ and if $\partial P\neq\varnothing$ then $P$ meets $\Pi$ orthogonally. Furthermore either $P$ is properly embedded or $P=\Pi$.  
        	 		As the ambient space is simply connected, we can always conclude that $P$ is two-sided and that this convergence happens with multiplicity one.     		
        	 	
        	 	
        	 	Towards a contradiction, suppose that either $P=\Pi$ or
        	 	$$\lambda_1(P\cap B^{\R^{n+1}}_{3/2}(0))\geqslant 0,$$ so that in particular $P$ is strictly stable on $B^{\R^{n+1}}_{4/3}(0)$. Notice that if $P=\Pi$ then the above is trivially true for variations in the whole of $\R^{n+1}$ (i.e. not just in $\Pi_1$). In either case, by smooth multiplicity one convergence we have that $\widetilde{P}_k$ is strictly stable on $B^{\R^{n+1}}_{4/3}(0)$ which implies (scaling back and using the second part of Remark \ref{rmk:coord}) that $P_k$ is stable on $B_{r_k + r_k^2}(p_k)$, a contradiction. In particular $P$ cannot be planar, and we have 
        	 	$$\lambda_1(P\cap B^{\R^{n+1}}_{3/2}(0))\leqslant -2\lambda_{\ast} < 0$$ for some $\lambda_{\ast}>0$, and hence for all $k$ sufficiently large
        	 	
        	 		$$\lambda_1(\widetilde{P}_k\cap B^{\R^{n+1}}_{3/2}(0))\leqslant -\lambda_{\ast}<0.$$ 
        	 		Hence a rescaling argument implies $$\lambda_1(P_k \cap (B_{2r_k}(p_k)))\longrightarrow -\infty.$$ 
        	 		
        	 		Obviously, the fact that for $n=2$ we must obtain a half-bubble relies on the half-space theorem \cite{HM90} (i.e. there cannot be a non-trivial bubble contained in an open half-space.)
        	 		
        	 		Finally, the argument for the last part is similar to the above except we always choose normal coordinates to carry out our blow-ups using Remark \ref{rmk:coord}. At the $r_k$ scale we have that $B_{\widehat{r}_k}(\widehat{p}_k)$ is similar to the balls 
        	 		$$B^{\R^{n+1}}_{\beta_k^{\pm 1}\widetilde{r}_k}(\widetilde{p}_k)$$ with $1\leqslant \widetilde{r}_k \leqslant C$ and $\widetilde{p}_k\in B^{\R^{n+1}}_C(0)$, at which point we can conclude that $\widetilde{r}_k \to \widetilde{r}\in [1,C]$ and $\widetilde{p}_k \to \widetilde{p}\in B^{\R^{n+1}}_C(0)$.    
        	 		Assuming that $(B^{\R^{n+1}}_{2\widetilde{r}}(\widetilde{p})\smallsetminus B^{\R^{n+1}}_2(0))\cap P$ is stable (or empty), we arrive at a contradiction in a similar fashion as above: first this implies that $B_{3\widetilde{r}/2}^{\R^{n+1}}(\widetilde{p})\smallsetminus B_2^{\R^{n+1}}(0)\cap P$ is strictly stable (or empty). But this domain contains the blown-up initial domain on which we are assuming that $P_k$ is unstable. Thus this must be non-empty, and by the smooth multiplicity one convergence we would obtain that $(B_{3\widetilde{r}/2}^{\R^{n+1}}(\widetilde{p})\smallsetminus B_2^{\R^{n+1}}(0))\cap \widetilde{P}_k$ is strictly stable, a contradiction. 
        	 		
        	 		Thus, we have proved that $$\lambda_1 ((B^{\R^{n+1}}_{3\widetilde{r}/2}(\widetilde{p})\smallsetminus B^{\R^{n+1}}_2(0))\cap P)\leqslant -2\lambda_{\ast}<0 $$ as before, for some $\lambda_{\ast}>0$. A rescaling argument concludes the final statement of the Theorem, once again appealing to Remark \ref{rmk:equivballs}.         
        	 \end{proof}
        	 	
        	 	
        	 	
        	 	
\begingroup\footnotesize
\begin{thebibliography}{48}
\bibitem[1]{ABCS18} L. Ambrozio, R. Buzano, A. Carlotto \& B. Sharp – “Geometric convergence results for closed minimal surfaces via bubbling analysis”, 2018, arXiv:1803.04956.
\bibitem[2]{ACS15} L. Ambrozio, A. Carlotto \& B. Sharp – “Compactness of the space of minimal hypersurfaces with bounded volume and $p$-th Jacobi eigenvalue”, J. Geom. Anal. 26 (2016), no. 4, p. 2591–2601.
\bibitem[3]{ACS17} \rule{2em}{.3pt} , “Compactness analysis for free boundary minimal hypersurfaces”, Calc. Var. Partial Differential Equations 57 (2018), no. 1, article ID 22 (39 pages).
\bibitem[4]{ACS18b} \rule{2em}{.3pt} , “Index estimates for free boundary minimal hypersurfaces”, Math. Ann. 370 (2018), no. 3-4, p. 1063–1078.
\bibitem[5]{BS17} R. Buzano \& B. Sharp – “Qualitative and quantitative estimates for minimal hypersurfaces with bounded index and area”, Trans. Amer. Math. Soc. 370 (2018), no. 6, p. 4373–4399.
\bibitem[6]{dCP79} M. do Carmo \& C. K. Peng – “Stable complete minimal surfaces in $\mathbf{R}^3$ are planes”, Bull. Amer. Math. Soc. (N.S.) 1 (1979), no. 6, p. 903–906.
\bibitem[7]{CFP14} J. Chen, A. Fraser \& C. Pang – “Minimal immersions of compact bordered Riemann surfaces with free boundary”, Trans. Amer. Math. Soc. 367 (2015), no. 4, p. 2487–2507.
\bibitem[8]{CT94} S. Y. Cheng \& J. Tysk – “Schrödinger operators and index bounds for minimal submanifolds”, Rocky Mountain J. Math. 24 (1994), no. 3, p. 977–996.
\bibitem[9]{CM14} O. Chodosh \& D. Maximo – “On the topology and index of minimal surfaces”, J. Differential Geom. 104 (2016), no. 3, p. 399–418.
\bibitem[10]{Cou40} R. Courant – “The existence of minimal surfaces of given topological structure under prescribed boundary conditions”, Acta Math. 72 (1940), p. 51–98.
\bibitem[11]{Cou50} R. Courant – Dirichlet’s principle, conformal mapping, and minimal surfaces, Springer-Verlag, New York-Heidelberg, 1977.
\bibitem[12]{DR16} C. De Lellis \& J. Ramic – “Min-max theory for minimal hypersurfaces with boundary”, Ann. Inst. Fourier (Grenoble) 68 (2018), no. 5, p. 1909–1986.
\bibitem[13]{EM08} N. Ejiri \& M. Micallef – “Comparison between second variation of area and second variation of energy of a minimal surface”, Adv. Calc. Var. 1 (2008), no. 3, p. 223–239.
\bibitem[14]{FC85} D. Fischer-Colbrie – “On complete minimal surfaces with finite Morse index in three-manifolds”, Invent. Math. 82 (1985), no. 1, p. 121–132.
\bibitem[15]{FCS80} D. Fischer-Colbrie \& R. Schoen – “The structure of complete stable minimal surfaces in $3$-manifolds of nonnegative scalar curvature”, Comm. Pure Appl. Math. 33 (1980), no. 2, p. 199–211.
\bibitem[16]{FPZ15} A. Folha, F. Pacard \& T. Zolotareva – “Free boundary minimal surfaces in the unit 3-ball”, Manuscripta Math. 154 (2017), no. 3-4, p. 359–409.
\bibitem[17]{FL14} A. Fraser \& M. M.-c. Li – “Compactness of the space of embedded minimal surfaces with free boundary in three-manifolds with nonnegative Ricci curvature and convex boundary”, J. Differential Geom. 96 (2014), no. 2, p. 183–200.
\bibitem[18]{FS16} A. Fraser \& R. Schoen – “Sharp eigenvalue bounds and minimal surfaces in the ball”, Invent. Math. 203 (2016), no. 3, p. 823–890.
\bibitem[19]{Fra00} A. M. Fraser – “On the free boundary variational problem for minimal disks”, Comm. Pure Appl. Math. 53 (2000), no. 8, p. 931–971.
\bibitem[20]{FGM16} B. Freidin, M. Gulian \& P. McGrath – “Free boundary minimal surfaces in the unit ball with low cohomogeneity”, Proc. Amer. Math. Soc. 145 (2017), no. 4, p. 1671–1683.
\bibitem[21]{GJ86A} M. Grüter \& J. Jost – “On embedded minimal disks in convex bodies”, Ann. Inst. H. Poincaré Anal. Non Linéaire 3 (1986), no. 5, p. 345–390.
\bibitem[22]{GZ18} Q. Guang \& X. Zhou – “Compactness and generic finiteness for free boundary minimal hypersurfaces”, 2018, arXiv:1803.01509.
\bibitem[23]{HM90} D. Hoffman \& W. H. Meeks, III – “The strong halfspace theorem for minimal surfaces”, Invent. Math. 101 (1990), no. 2, p. 373–377.
\bibitem[24]{JM83} L. P. Jorge \& W. H. Meeks, III – “The topology of complete minimal surfaces of finite total Gaussian curvature”, Topology 22 (1983), no. 2, p. 203–221.
\bibitem[25]{Jo86} J. Jost – “Existence results for embedded minimal surfaces of controlled topological type. I”, Ann. Scuola Norm. Sup. Pisa Cl. Sci. (4) 13 (1986), no. 1, p. 15–50.
\bibitem[26]{Jo86b} \rule{2em}{.3pt} , “Existence results for embedded minimal surfaces of controlled topological type. II”, Ann. Scuola Norm. Sup. Pisa Cl. Sci. (4) 13 (1986), no. 3, p. 401–426.
\bibitem[27]{Jo87} \rule{2em}{.3pt} , “Existence results for embedded minimal surfaces of controlled topological type. III”, Ann. Scuola Norm. Sup. Pisa Cl. Sci. (4) 14 (1987), no. 1, p. 165–167.
\bibitem[28]{KL17} N. Kapouleas \& N. Li – “Free boundary minimal surfaces in the unit three-ball via desingularization of the critical catenoid and equatorial disk”, 2017, arXiv:1709.08556.
\bibitem[29]{KW17} N. Kapouleas \& D. Wygul – “Free-boundary minimal surfaces with connected boundary in the 3-ball by tripling the equatorial disc”, 2017, arXiv:1711.00818.
\bibitem[30]{Ket16A} D. Ketover – “Free boundary minimal surfaces of unbounded genus”, 2016, arXiv:1612.08691.
\bibitem[31]{Li15} M. M.-c. Li – “A general existence theorem for embedded minimal surfaces with free boundary”, Comm. Pure Appl. Math. 68 (2015), no. 2, p. 286–331.
\bibitem[32]{LZ16B} N. Li \& X. Zhou – “Min-max theory for free boundary minimal hypersurfaces I – regularity theory”, 2016, arXiv:1611.02612.
\bibitem[33]{Lim17} V. Lima – “Bounds for the Morse index of free boundary minimal surfaces”, 2017, arXiv: 1710.10971.
\bibitem[34]{LR89} F. J. López \& A. Ros – “Complete minimal surfaces with index one and stable constant mean curvature surfaces”, Comment. Math. Helv. 64 (1989), no. 1, p. 34–43.
\bibitem[35]{LR91} \rule{2em}{.3pt} , “On embedded complete minimal surfaces of genus zero”, J. Differential Geom. 33 (1991), no. 1, p. 293–300.
\bibitem[36]{MNS13} D. Maximo, I. Nunes \& G. Smith – “Free boundary minimal annuli in convex three-manifolds”, J. Differential Geom. 106 (2017), no. 1, p. 139–186.
\bibitem[37]{Oss63} R. Osserman – “On complete minimal surfaces”, Arch. Rational Mech. Anal. 13 (1963), p. 392– 404.
\bibitem[38]{Oss64} \rule{2em}{.3pt} , “Global properties of minimal surfaces in $E^3$ and $E^n$”, Ann. of Math. (2) 80 (1964), p. 340–364.
\bibitem[39]{Pog81} A. V. Pogorelov – “On the stability of minimal surfaces in Lobachevski˘ı space”, Dokl. Akad. Nauk 354 (1997), no. 6, p. 742–744.
\bibitem[40]{RV95} A. Ros \& E. Vergasta – “Stability for hypersurfaces of constant mean curvature with free boundary”, Geom. Dedicata 56 (1995), no. 1, p. 19–33.
\bibitem[41]{SS81} R. Schoen \& L. Simon – “Regularity of stable minimal hypersurfaces”, Comm. Pure Appl. Math. 34 (1981), no. 6, p. 741–797.
\bibitem[42]{Sc83} R. M. Schoen – “Uniqueness, symmetry, and embeddedness of minimal surfaces”, J. Differential Geom. 18 (1983), no. 4, p. 791–809 (1984).
\bibitem[43]{Sha15} B. Sharp – “Compactness of minimal hypersurfaces with bounded index”, J. Differential Geom. 106 (2017), no. 2, p. 317–339.
\bibitem[44]{Str84} M. Struwe – “On a free boundary problem for minimal surfaces”, Invent. Math. 75 (1984), no. 3, p. 547–560.
\bibitem[45]{T89} J. Tysk – “Finiteness of index and total scalar curvature for minimal hypersurfaces”, Proc. Amer. Math. Soc. 105 (1989), no. 2, p. 429–435.
\bibitem[46]{Wan99} G. Wang – “Birkhoffminimax principle for minimal surfaces with a free boundary”, Math. Ann. 314 (1999), no. 1, p. 89–107.
\bibitem[47]{Whi09} B. White – “Which ambient spaces admit isoperimetric inequalities for submanifolds?”, J. Differential Geom. 83 (2009), no. 1, p. 213–228.
\bibitem[48]{W15} \rule{2em}{.3pt} , “On the compactness theorem for embedded minimal surfaces in 3-manifolds with locally bounded area and genus”, Comm. Anal. Geom. 26 (2018), no. 3, p. 659–678.
\end{thebibliography}
\endgroup

\end{document}
